\documentclass[a4paper,11pt]{article}
\usepackage[utf8]{inputenc}
\usepackage[T1]{fontenc}
\usepackage[english]{babel}
\usepackage[left=2.5cm,right=2.5cm,top=2cm,bottom=2cm]{geometry}
\usepackage{hyperref}
\usepackage{xcolor}
\usepackage{microtype}
\usepackage{lmodern}

\definecolor{primary}{RGB}{0, 129, 145}
\hypersetup{colorlinks=true, urlcolor=primary}

\begin{document}
\noindent \textbf{Lo\"ic Aron MBASSI EWOLO} \\
ENSEA -- Double Dipl\^ome Ing\'enieur \\
\href{mailto:loicaron.mbassiewolo@ensea.fr}{loicaron.mbassiewolo@ensea.fr} \quad (+33) 7 59 52 33 98 \\
\href{https://github.com/Nameless0l}{github.com/Nameless0l} \quad \href{https://mbassiloic.tech}{mbassiloic.tech}

\vspace{1em}\noindent Saint-Maurice, le \today
\vspace{1em}\noindent \textbf{Objet : Candidature Alternance -- Ing\'enieur Avant-Vente IoT \& Smart Metering}
\vspace{1.5em}\noindent Madame, Monsieur,

\vspace{0.8em}\noindent J'ai d\'eploy\'e des mod\`eles TensorFlow Lite sur microcontr\^oleurs (Arduino, Raspberry Pi) avec acquisition et traitement de donn\'ees capteurs IoT en temps r\'eel dans le cadre du Stage IoT Lab ENSPY, et int\'egr\'e un syst\`eme complet de capteurs IMU/Flex sur STM32 (protocoles I2C/SPI) dans le projet Harmony Gloves au Labo IOTA ENSEA. Formation double dipl\^ome ENSEA/Polytechnique Yaound\'e, sp\'ecialit\'e syst\`emes embarqu\'es et math\'ematiques appliqu\'ees. Actuellement en stage INRIA Saclay (projet PRIMO) sur pipelines avanc\'es et analyse statique.

\vspace{0.8em}\noindent Veolia, acteur mondial de la gestion des ressources naturelles, d\'eploie des r\'eseaux IoT et de smart metering pour optimiser la distribution d'eau et d'\'energie \`a grande \'echelle. Cette alternance en avant-vente repr\'esente une opportunit\'e unique d'appliquer ma maîtrise des architectures IoT compl\`etes et des protocoles LPWAN (LoRa, Sigfox) \`a la conception de solutions de mesure intelligente pour services publics.

\vspace{0.8em}\noindent Mon exp\'erience Urban Waste Management (1\textsuperscript{er} prix CITS National hackathon, 75 \'equipes) a valid\'e ma capacit\'e \`a concevoir des architectures IoT bout-en-bout : capteurs capacitifs, r\'eseaux de communication, backend centralis\'e, mod\`eles ML pour pr\'ediction et d\'etection anomalies. J'ai \'egalement optimis\'e des mod\`eles ML pour contraintes hardware dans l'IoT embarqu\'e, avec d\'eploiement temps r\'eel sur cibles h\'et\'erog\`enes. En FPGA, j'ai impl\'ement\'e des syst\`emes SoPC complets (Nios V, I2C), une exp\'erience transposable aux syst\`emes smart metering.

\vspace{0.8em}\noindent Je suis disponible en alternance d\`es septembre 2026 pour une dur\'ee de 12 mois, mobile en Île-de-France. Je serais ravi de discuter de comment mes comp\'etences en IoT embarqu\'e, protocoles LPWAN et architectures de smart metering peuvent contribuer aux solutions innovantes de Veolia pour la transition \'ecologique.

\vspace{1.5em}\noindent Veuillez agr\'eer, Madame, Monsieur, l'expression de mes salutations distingu\'ees.
\vspace{2em}\noindent \textbf{Lo\"ic Aron MBASSI EWOLO}
\end{document}
